\documentclass{article}
\usepackage[T1]{fontenc}
\usepackage{lmodern}
\usepackage{graphicx}
\usepackage{amsmath,amssymb}
\usepackage[a4paper,margin=2cm]{geometry}
\usepackage[absolute,overlay]{textpos}
\setlength{\TPHorizModule}{1mm}
\setlength{\TPVertModule}{1mm}
\usepackage[indonesian]{babel}
\usepackage{hyperref}
\setlength{\emergencystretch}{2em}
% unit: Halaman 18

\begin{document}

% === Halaman 18 (python/native) ===

18

Akhirnya, campurkan L dan S sebagai berikut: 

i  0 

j  0 

X  0 

Y  0 

n  3 * max(r, c) 

for k  1 to n do 

    S[i]  (S[i] + X + Y) <<< 3 

    X  S[i] 

    i  (1 + 1) mod (t) 

    L[j]  (L[j] + X + Y) <<< ( X + Y) 

    Y  L[j] 

    j  (j + 1) mod (c) 

endfor

\end{document}
